% Section Introduction
Computational models and simulations (CM\&S) are increasingly used across the medical device lifecycle, from early stages in design to optimization of in silico testing and post-market surveillance \cite{morrisonFrontiers}. These tools offer cost-effective, flexible, and scalable alternatives to traditional experimentation. However, their use in high risk regulatory decision making remains limited \cite{fda_credibility_program, fda2022Success}. To be accepted by regulatory agencies such as the U.S. Food and Drug Administration (FDA), computational models must be supported by rigorous, transparent evidence demonstrating that they are appropriate for their use case and capable of producing trustworthy results. 

Establishing this trust is particularly challenging for systems involving multiphysics interactions, coupled domains, and inherent uncertainties, especially when experimental data are sparse, costly, or difficult to obtain. In such settings, verification, validation, and uncertainty quantification (VVUQ) are essential. However, the required level of rigor must be aligned with the model’s intended use and associated risk. For medium to low risk applications, comprehensive experimental studies or high-fidelity simulations may be impractical. In these cases, credibility depends not only on numerical rigor but also on transparency, including clearly communicating assumptions, limitations, and modeling decisions to support risk-informed interpretation.

To provide context for the level of verification performed, this study adopts a risk-informed perspective consistent with the ASME V\&V 40 framework. For the present application, the model is characterized by a low decision consequence and moderate model influence, corresponding to a medium–low risk classification. Under this classification, a moderate to minimal level of verification would typically be sufficient. However, in this work, additional rigor is pursued where computationally feasible. In particular, simplified geometries and isolated physical configurations are used to enable more detailed assessment of numerical behavior, including estimation of discretization error and evaluation of observed order of accuracy. This approach provides a level of verification beyond minimal while remaining practical.

This study focuses specifically on verification as a foundational component of the VVUQ process, with validation and uncertainty quantification addressed in a subsequent companion study. In particular, this study presents a structured and transparent verification study of a computational model for a medical device application involving heated internal flow. The system is modeled as a conjugate heat transfer (CHT) problem, where laminar fluid flows through a heated solid casing and exchanges thermal energy before exiting the device. Simulations are performed using the open-source finite volume solver \texttt{OpenFOAM}, specifically the \texttt{chtMultiRegionFoam} solver for coupled fluid-solid heat transfer  \cite{openfoamSolvers, chtMultiRegion}.

The verification approach focuses on both code and solution verification to establish numerical correctness and quantify discretization error. This includes the use of analytical problems, systematic mesh and timestep refinement studies, and regression-based error estimation. Particular emphasis is placed on transient, coupled configurations, where spatial and temporal resolution interact and cannot be assessed independently.

This work is grounded in the ASME V\&V 40 framework for credibility assessment in medical device applications and aligns with FDA-recognized standards, including V\&V 10, V\&V 20, and V\&V 40 \cite{vv10, vv20, vv40, pathmanathanAA, vv10_fdarecognized, vv20_fdarecognized, vv40_fdarecognized, fda_guidance_2023}. While these standards define high-level credibility principles, they do not prescribe specific verification workflows. To address this gap, the methodology presented here draws on best practices from high-consequence fields such as aerospace and nuclear engineering, where VVUQ processes are well established due to similar experimental and modeling constraints \cite{OberkampfRoy2013, Roache2009, AIAA-G-077-1998}. Additional details on the risk-informed credibility assessment, including V\&V 40 gradation selection and justification, are provided in publicly accessible repository materials \cite{edds_credibility_repo}.


By openly documenting simulation workflows, input files, and post-processing procedures, this study provides a reproducible case study in verification for medical device modeling. The goal is to demonstrate how rigorous, transparent, and risk-informed verification can be achieved within realistic computational and resource constraints, supporting broader adoption of CM\&S in regulatory and engineering applications.